\documentclass[10pt]{article}
\usepackage[T1]{fontenc}
\usepackage[utf8]{inputenc}
\usepackage[french]{babel}
\usepackage{amsmath,amssymb,mathrsfs,mathpazo}
\usepackage{geometry}
\geometry{a4paper,margin=2cm}
\usepackage{graphicx,array,tabularx,tabularray,colortbl,multirow,multicol,booktabs,enumitem,pifont,ulem,fancybox,diagbox,cancel}
\usepackage{pstricks,pstricks-add,pst-plot,pst-tree,pst-node,pst-text,pst-func,pst-math,pst-bspline,pst-3dplot,pst-eucl,pst-solides3d}
\let\pscalculate\undefined
\usepackage{pst-bezier}
\usepackage{tikz,pgfplots,tkz-tab}
\usetikzlibrary{arrows,arrows.meta,patterns,calc,positioning,trees,decorations.markings,shapes.geometric,shapes.misc,shapes.arrows,intersections,decorations.pathmorphing,decorations.pathreplacing,backgrounds,shadows,babel}
\pgfplotsset{compat=1.18}
\usepgfplotslibrary{fillbetween}
\usepackage[np]{numprint}
\usepackage{listings,makecell,pas-tableur,twoopt}
\let\barre\undefined
\usepackage[pstricks]{tablvar}
\definecolor{darkgreen}{RGB}{0,100,0}
\providecommand{\red}{\color{red}}
\providecommand{\blue}{\color{blue}}
\providecommand{\green}{\color{green}}
\providecommand{\black}{\color{black}}
\providecommand{\white}{\color{white}}
\providecommand{\gray}{\color{gray}}
\providecommand{\cyan}{\color{cyan}}
\providecommand{\magenta}{\color{magenta}}
\providecommand{\yellow}{\color{yellow}}
\providecommand{\violet}{\color{violet}}
\providecommand{\orange}{\color{orange}}
\providecommand{\euro}{\texteuro}
\providecommand{\EUR}{\texteuro}
\providecommand{\ieme}{\textsuperscript{e}}
\providecommand{\degres}{\ensuremath{^\circ}}
\providecommand{\degree}{\ensuremath{^\circ}}
\providecommand{\vv}[1]{\overrightarrow{#1}}
\providecommand{\Coord}[1]{#1}
\providecommand{\decosix}{}
\providecommand{\starredbullet}{\ensuremath{\bullet}}
\pagestyle{empty}
\setlength{\parindent}{0pt}
\usepackage{adjustbox}\DeclareUnicodeCharacter{2236}{:}\DeclareUnicodeCharacter{2190}{\ensuremath{\leftarrow}}
\providecommand{\R}{\mathbb{R}}
\providecommand{\N}{\mathbb{N}}
\providecommand{\Z}{\mathbb{Z}}
\providecommand{\Q}{\mathbb{Q}}
\providecommand{\D}{\mathbb{D}}
\providecommand{\C}{\mathbb{C}}
\providecommand{\e}{{\mathrm{e}}}
\providecommand{\vect}[1]{\overrightarrow{#1}}
\providecommand{\vectt}[1]{\overrightarrow{\mathrm{#1}}}
\providecommand{\vvt}[1]{\overrightarrow{\mathrm{#1}}}
\providecommand{\np}[2][]{{#2~#1}}
\providecommand{\vv}[1]{{\overrightarrow{#1}}}
\renewcommand{\d}{{\mathrm{d}}}
\providecommand{\degres}{{^\circ}}
\providecommand{\degre}{{^\circ}}
\providecommand{\blue}{\color{blue}}
\providecommand{\red}{\color{red}}
\providecommand{\enskip}{\quad}
\providecommand{\llbracket}{\lbrack\!\lbrack}
\providecommand{\rrbracket}{\rbrack\!\rbrack}
\providecommand{\euro}{{\text{€}}}
\providecommand{\no}{{\text{n°}}}
\providecommand{\interval}[2]{{[#1;#2]}}
\providecommand{\intervFO}[2]{{[#1;#2[}}
\providecommand{\dotfill}{\ldots}

\providecommand{\ds}{\displaystyle}
\providecommand{\textding}[1]{#1}
\providecommand{\Oij}{\(\mathrm{O};\vec{\imath},\vec{\jmath}\)}
\providecommand{\Oijk}{\(\mathrm{O};\vec{\imath},\vec{\jmath},\vec{k}\)}
\providecommand{\Ouv}{\(\mathrm{O};\vec{u},\vec{v}\)}
\providecommand{\intervFF}[2]{\ensuremath{[#1;#2]}}
\renewcommand{\euro}{\ifmmode\text{\texteuro}\else\texteuro\fi}
\DeclareUnicodeCharacter{00B0}{\ifmmode^{\circ}\else\textdegree\fi}
\DeclareUnicodeCharacter{20AC}{\ifmmode\text{\texteuro}\else\texteuro\fi}
% Énoncés : Éducation nationale. Transcription et corrigés : APMEP, Denis Vergès et contributeurs. Découpage et classement : CoopMaths, Sébastien Lozano, Éric Elter et contributeurs. Adaptation de la présentation : Maths Explorer.
% CC BY-SA : https://coopmaths.fr/www/about/
% https://coopmaths.fr/alea/static/eam/2025/tex/eam_specifique_2025_06_sujet0v1_3_cor.tex
\begin{document}
\noindent\textit{APMEP / CoopMaths. CC BY-SA. Présentation adaptée par Maths Explorer.}\par\medskip
\medskip

Un vendeur de voitures possède un stock de \np{1000} voitures dont les caractéristiques sont résumées dans le tableau ci-dessous.

\medskip

\begin{tabularx}{\linewidth}{|l|*{4}{>{\centering \arraybackslash}X|}}\hline
&Blanche &Noire &Rouge&TOTAL\\ \hline
Française &150 &$x$&400&750\\ \hline
Étrangère &100& 50 &100& 250\\ \hline
TOTAL &250&250 &500&\np{1000}\\ \hline
\end{tabularx}

\begin{enumerate}
\item 
Le nombre $x$ représente le nombre de voitures françaises noires.

Il y a 250 voitures noires, dont 50 étrangères; donc $x=250-50=200$.

\item 
Il y a \np{1000} voitures dans le stock dont 250 voitures noires.

Le pourcentage de voitures noires parmi les voitures du stock est donc égal à\\
 $\dfrac{250}{\np{1000}}$ c'est-à-dire $\dfrac{25}{100}$ soit 25\,\%.

\item 
Il y a \np{1000} voitures dans le stock dont 50 voitures noires étrangères.

Le pourcentage de voitures noires étrangères parmi les voitures du stock est donc égal à
 $\dfrac{50}{\np{1000}}$ c'est-à-dire $\dfrac{5}{100}$ soit 5\,\%.

\item 
Il y a 750 voitures françaises dont 150 voitures blanches.

Le pourcentage de voitures blanches parmi les voitures françaises est donc égal à\\
 $\dfrac{150}{\np{750}}$ c'est-à-dire $\dfrac{1}{5}$ ou $\dfrac{20}{100}$ soit 20\,\%.

\item 
Il y a 250 voitures blanches dont 150 voitures françaises.

Le pourcentage de voitures françaises parmi les voitures blanches  est donc égal à\\
 $\dfrac{150}{\np{250}}$ c'est-à-dire $\dfrac{15}{25}$ ou $\dfrac{60}{100}$ soit 60\,\%.

\item 
\begin{list}{\textbullet}{ Alice et Benoît jouent au jeu suivant:}
\item Alice choisit au hasard une voiture parmi les voitures Françaises. \\
Elle remporte 1 euro si ce n'est pas une voiture rouge.

Parmi les 750 voitures françaises, il y en a $750-400=350$ qui ne sont pas rouges. La probabilité pour Alice de gagner est donc $\dfrac{350}{750}$.

\item Benoit choisit au hasard une voiture parmi les voitures Blanches. \\
Il remporte 1 euro si c'est une voiture étrangère.

Parmi les 250 voitures blanches, il y en a 100 qui sont étrangères. La probabilité pour Benoit de gagner est donc $\dfrac{100}{250}$.
\end{list}

Or $\dfrac{100}{250} = \dfrac{300}{750} < \dfrac{350}{750}$,
donc c'est Alice qui a le plus de chance de remporter 1 euro.
\end{enumerate}

\bigskip

\end{document}
